\section{Gemini 4 Neonの利用評価が活発化}
Googleが発表したGemini 4 Neonを巡り、観測窓内に多数の直接言及があった。CarbonとHydrogenの間に位置づけられ、100万トークンのコンテキストと4段階のthinking level、コーディングエージェント向けの性能が特徴とされる。利用者投稿ではArtificial AnalysisのIntelligence Indexで61点・3位という評価や、100万トークン当たり入力3ドル・出力15ドルという価格が話題になった。順位や使用感は投稿者の報告として扱う。公式発表は観測窓開始約25分前のため、窓外である。
\dailyxsource{Google AI公式投稿（観測窓開始直前）}{https://x.com/GoogleAI/status/2123168242797809944}

\section{AnthropicのIPO価格と取引開始を巡る投稿}
AnthropicのIPO価格が1株83ドルに決まり、Nasdaqで銘柄ANTとして取引を開始したとする投稿が複数見られた。初日の値動きや調達額に関する言及もあり、生成AI企業の資本市場での評価に関心が集まっている。ただしGrokは直接の公式X発表を特定しておらず、Bloombergなどの報道に言及する二次投稿が中心である。価格、銘柄、上場状況はいずれも観測された主張であり、独立確認済みの事実とは区別する。
\dailyxsource{IPOに関する窓内の投稿例}{https://x.com/kakatatayaya/status/2123538790912606237}

\section{Runwayがロボット制御向けPraxis-1を公開}
Runwayの開発者アカウントがロボティクス向けworld action model「Praxis-1」を紹介した。27BパラメータのMoE diffusion transformerとして、動画とアクションを同時に予測する設計が説明されている。GitHubでのオープンウェイト公開や家庭用ロボットのデモも話題となった。NVIDIA RoboLab-120における性能についての言及は、投稿者側の主張として扱う。動画生成とロボット制御を結びつける研究の進展が注目点である。
\dailyxsource{Runway Developers公式投稿}{https://x.com/RunwayMLDevs/status/2123516008950366531}

\section{Latent Lattice Language Models論文が話題}
Latent Lattice Language Models（LLLMs）に関する研究紹介が観測された。各Transformer層に持続的な隠れ状態を加えることで、モデルが内部情報を保持する仕組みを拡張するという。パラメータ効率やperplexityの改善を示す主張も紹介されたが、性能の詳細は著者や紹介投稿による説明の範囲にとどまる。Grokの収集結果には正確なarXiv IDや特定可能な一次X投稿URLがないため、本号ではリンクを補わない。

\section{DaisyRose-27Bのエージェント性能を巡る報告}
RadixArkのDaisyRose-27Bについて、Qwen系モデルを基盤としてエージェントタスク向けに強化学習を実施したとの投稿が見られた。約1,200タスクでの学習やTerminal-Benchの改善が主張され、オープンモデルを実務的なエージェントへ最適化する方向性が注目されている。ただし公式の一次投稿やモデル配布先はGrokの観測では特定されていない。訓練規模や性能改善は現段階では投稿上の情報として扱う。

\section{小型音声認識モデルRumble-2-Flash}
Reson8のRumble-2-Flashを紹介する投稿が観測された。約5億パラメータのオープンウェイト音声認識モデルとして、Apache 2.0での公開、40言語以上への対応、実環境音声での性能が主張されている。軽量な多言語音声認識はローカル処理や低遅延アプリケーションでの利用可能性がある一方、今回の収集では公式モデルカードやベンチマークの一次資料を確認できていない。数値やライセンスの扱いには留意が必要である。

\section{OpenAIが研究でのAI利用を報告}
OpenAIの公式アカウントは、自社研究におけるAI利用の増加について投稿した。エージェントが長時間のタスクを完了する能力や、安全性に関する人間側の取り組みの増加に言及し、再帰的自己改善（RSI）を測る指標の公開を提唱したという。研究活動そのものをAIが支援する流れが主題であり、能力の評価方法と安全性の測定をどう結びつけるかが論点となる。内容はOpenAI自身の報告に基づく。
\dailyxsource{OpenAI公式投稿}{https://x.com/OpenAI/status/2123191365983248653}

\section{LLMとタンパク質構造予測の関連研究}
LLMの内部状態とタンパク質の折りたたみを結びつける研究論文について、X上で紹介投稿が見られた。TM-scoreなどの評価指標への言及もあったが、Grokは論文の一次資料や具体的な投稿URLを確認していない。どのような入力から構造を予測するのか、既存手法との比較が成立するのかについては判断できないため、現時点では未確認の研究紹介として記録する。